\section{Limitations and Research Implications}

\subsection{What the evidence does not establish}

The sample begins in 2016 and yields 66 evaluable months because five years are consumed
by state formation. Fourteen primary-proxy Q4 observations provide limited power for a
tail-state claim. The study uses daily data and monthly sleeves; it says nothing direct
about minute-level or order-book reversal. It also tests one published momentum window,
one reversal window, one quantile boundary, and two market proxies. Failure to replicate
this rule is not proof that every volatility-aware allocation is ineffective.

Nor does the study show that Chinese futures lack momentum. The frozen momentum sleeve is
positive in the sample. It does not rule out reversal at other horizons or during
particular market events. Finally, the cross-sectional anatomy diagnostics are
descriptive. They show that Q4 is more systemic, but they do not causally identify which
economic shock produced each state.

\subsection{What can be learned from a failed replication}

The failure changes the object of research. The useful question is no longer ``which
volatility percentile maximizes the backtest?'' It is ``which high-volatility states
destroy trends, and which create trends?'' That question is both more economic and more
falsifiable.

A future state model should be low-dimensional and interpretable. It might combine shock
intensity with pre-specified measures of market direction, breadth, and shock stage. Each
additional dimension must first explain the future reversal-minus-momentum spread across
time partitions and reasonable market proxies. Only a mechanism that passes those gates
should be translated into a router and tested on an untouched period.

This sequence matters for inference. Adding indicators until a router improves on the
same 66 months would merely transfer overfitting from a scalar threshold to a state
classifier. Mechanism diagnostics, portfolio selection, and final validation must remain
separate stages.

\subsection{General implication for model transportability}

State variables are not portable solely because their mathematical definitions are
portable. A volatility standard deviation has the same formula in equities and futures,
but its economic content depends on the aggregation object and on which shocks dominate
the market. Cross-market replication should therefore audit three layers:

\begin{enumerate}
  \item \textbf{Measurement invariance:} does the state have a unique or stable market
  definition?
  \item \textbf{Mechanism invariance:} does the state proxy the same economic forces?
  \item \textbf{Decision invariance:} does crossing the state boundary select the same
  return-generating behavior?
\end{enumerate}

The present replication fails all three strongly enough that endpoint outperformance in
one proxy is not persuasive evidence of transportability.
